\documentclass[11pt]{article}
\usepackage[T1]{fontenc}
\usepackage[utf8]{inputenc}
\usepackage{lmodern}
\usepackage{geometry}
\usepackage{hyperref}
\usepackage{amsmath}
\usepackage{amssymb}
\usepackage{enumitem}
\geometry{margin=1in}

\title{Implementation of the LDG Method in Exasim}
\author{}
\date{}

\begin{document}
\maketitle

\section*{1. Summary}

Exasim implements the LDG method by rewriting the PDE into a first-order
system, recovering the auxiliary gradient variable $q$ locally, then
assembling the $u$ residual from element-volume and face-flux terms. The
core LDG path is in:

\begin{itemize}[leftmargin=*]
\item \texttt{backend/Discretization/residual.cpp}
\item \texttt{backend/Discretization/qresidual.cpp}
\item \texttt{backend/Discretization/uresidual.cpp}
\item \texttt{backend/Discretization/getuhat.cpp}
\item \texttt{backend/Discretization/geometry.cpp}
\end{itemize}

The PDE-specific terms in the weak formulation are injected through the
model-driver layer in:

\begin{itemize}[leftmargin=*]
\item \texttt{backend/Model/ModelDispatch/model\_drivers\_abi.cpp}
\end{itemize}

\section*{2. Assumptions}

\begin{itemize}[leftmargin=*]
\item I am describing the LDG residual implementation, not the HDG global
trace-system implementation.
\item The PDE is written in first-order form with unknowns $u$ and $q$, where
$q$ represents gradients or gradient-like auxiliary variables.
\item I use the Exasim code sign conventions as implemented in the residual
assembly. These are equivalent to standard LDG weak forms, but the algebraic
residual may be written with the opposite overall sign from a textbook
presentation.
\item Arrays are stored in Exasim's flattened column-major style and are
evaluated by:
\begin{itemize}[leftmargin=*]
\item gather nodal data
\item interpolate to Gauss points
\item evaluate model terms
\item apply metric/Jacobian factors
\item project back to nodal residuals
\end{itemize}
\item $w$ in Exasim may also denote an additional algebraic/auxiliary field
\texttt{wdg}; that is separate from the usual test function notation. Below I
will use $v$ and $\phi$ for test functions.
\end{itemize}

\section*{3. Design reasoning}

A standard LDG formulation for a diffusion-type or mixed first-order PDE is:

\begin{align*}
q - \nabla u &= 0, \\
u_t + \nabla \cdot F(u,q,o,w) &= S(u,q,o,w).
\end{align*}

A typical elementwise weak form is:

\begin{itemize}[leftmargin=*]
\item $q$ equation:
\[
(q, v)_K + (u, \nabla \cdot v)_K - \langle \hat{u}, v \cdot n \rangle_{\partial K} = 0
\]
\item $u$ equation:
\[
(u_t, \phi)_K + (S, \phi)_K + (F, \nabla \phi)_K - \langle \hat{f}, \phi \rangle_{\partial K} = 0
\]
\end{itemize}

Exasim implements each of these terms explicitly.

\subsection*{A. High-level LDG execution order}

For the $u$ residual, Exasim does the following in
\texttt{backend/Discretization/residual.cpp}:

\begin{enumerate}[leftmargin=*]
\item \texttt{GetUhat(...)}
\item \texttt{GetQ(...)} if \texttt{common.ncq > 0}
\item \texttt{GetW(...)} if \texttt{common.ncw > 0}
\item \texttt{RuElem(...)}
\item \texttt{RuFace(...)}
\end{enumerate}

So the LDG structure is:

\begin{itemize}[leftmargin=*]
\item build face trace/boundary states
\item recover local gradient variable $q$
\item optionally recover algebraic field $w$
\item assemble the $u$ residual
\end{itemize}

That is the core LDG implementation pattern in Exasim.

\subsection*{B. Implementation of the weak $q$ equation}

This is handled in \texttt{backend/Discretization/qresidual.cpp}.

\paragraph{Term 1: $(u, \nabla \cdot v)_K$}

Implemented in \texttt{RqElemBlock(...)}.

What the code does:

\begin{itemize}[leftmargin=*]
\item gathers element nodal $u$ from \texttt{sol.udg}
\item interpolates $u$ from nodal points to Gauss points with
\texttt{Node2Gauss(...)}
\item applies geometric transformation with \texttt{ApplyXx3(...)}
\item projects back to element nodes with \texttt{Gauss2Node(...)}
\end{itemize}

Interpretation:

\begin{itemize}[leftmargin=*]
\item \texttt{Node2Gauss(...)} evaluates $u_h$ at quadrature points
\item \texttt{ApplyXx3(...)} applies the mapped derivative operators through
the metric tensor \texttt{Xx}
\item \texttt{Gauss2Node(...)} performs the quadrature-weighted test-function
contraction back to the discrete residual basis
\end{itemize}

So \texttt{RqElemBlock(...)} is the implementation of the volume term
$(u, \nabla \cdot v)_K$.

\paragraph{Term 2: $\langle \hat{u}, v \cdot n \rangle_{\partial K}$}

Implemented in \texttt{RqFaceBlock(...)}.

What the code does:

\begin{itemize}[leftmargin=*]
\item gathers face values of \texttt{uh} with \texttt{GetElemNodes(...)}
\item interpolates face \texttt{uh} to face Gauss points with
\texttt{Node2Gauss(...)}
\item applies \texttt{jac * n} with \texttt{ApplyJacNormal(...)}
\item projects back to the nodal residual with \texttt{Gauss2Node(...)}
\end{itemize}

This is exactly the discrete face integral of the LDG gradient equation.

\paragraph{Term 3: $(q, v)_K$}

Exasim does not assemble this as a separate explicit mass-matrix term in
\texttt{RqElemBlock}. Instead, after building the right-hand side
contributions from the volume and face terms, it applies the inverse mass
matrix in \texttt{GetQ(...)} inside
\texttt{backend/Discretization/residual.cpp}:

\begin{itemize}[leftmargin=*]
\item face contributions are scattered into the element residual by
\texttt{PutFaceNodes(...)}
\item \texttt{ApplyMinv(...)} applies the local inverse mass matrix
\end{itemize}

This is the discrete realization of solving:

\[
M q = RHS
\]

for each element, where $M$ is the element mass matrix from $(q, v)_K$.

So in Exasim:

\begin{itemize}[leftmargin=*]
\item $(q, v)_K$ is implemented through the local mass inverse
\texttt{ApplyMinv(...)}
\item $(u, \nabla \cdot v)_K$ is implemented by \texttt{RqElemBlock(...)}
\item $\langle \hat{u}, v \cdot n \rangle_{\partial K}$ is implemented by
\texttt{RqFaceBlock(...)}
\end{itemize}

After that, the recovered $q$ is inserted into \texttt{sol.udg} with
\texttt{ArrayInsert(...)}.

Important structural point:

\begin{itemize}[leftmargin=*]
\item in LDG, $q$ is local and recovered elementwise
\item it is not a globally coupled trace unknown
\end{itemize}

\subsection*{C. Construction of $\hat{u}$ in LDG}

This is handled in \texttt{backend/Discretization/getuhat.cpp}.

\texttt{GetUhat(...)} builds the face state \texttt{sol.uh}.

\paragraph{Interior faces}

For \texttt{ib == 0}, the code copies interior face data from
\texttt{sol.udg} into \texttt{sol.uh}.

Interpretation:

\begin{itemize}[leftmargin=*]
\item LDG in Exasim does not solve a global trace equation for $\hat{u}$
\item $\hat{u}$ is just the face state passed into the numerical flux machinery
\end{itemize}

\paragraph{Boundary faces}

For \texttt{ib != 0}, the code:

\begin{itemize}[leftmargin=*]
\item computes face geometry with \texttt{FaceGeom1D/2D/3D(...)}
\item gathers boundary-side \texttt{udg} and \texttt{odg}
\item calls \texttt{UbouDriver(...)}
\item writes the boundary trace into \texttt{sol.uh}
\end{itemize}

So the boundary trace $\hat{u}$ is entirely supplied by the user/model
boundary operator.

This is the implementation of the boundary data entering the LDG weak form.

\subsection*{D. Implementation of the weak $u$ equation}

This is handled in \texttt{backend/Discretization/uresidual.cpp}.

\paragraph{Term 1: time term $(u_t, \phi)_K$}

In \texttt{RuElemBlock(...)}, for time-dependent problems Exasim forms a
time-source-like contribution using \texttt{sol.sdg} and the time
discretization factor:

\begin{itemize}[leftmargin=*]
\item \texttt{ArrayAXPBY(...)} forms something equivalent to
\texttt{sdg - dtfactor * u}
\item if \texttt{common.tdfunc == 1}, \texttt{TdfuncDriver(...)} modifies this
term
\end{itemize}

So the temporal weak term is implemented as part of the element source
assembly rather than as a separate conceptual block.

\paragraph{Term 2: source term $(S, \phi)_K$}

Also in \texttt{RuElemBlock(...)}:

\begin{itemize}[leftmargin=*]
\item \texttt{SourceDriver(...)} evaluates the PDE source at Gauss points
\end{itemize}

This source is then combined into the element integrand before projection back
to nodal residuals.

So $(S, \phi)_K$ is implemented by:

\begin{itemize}[leftmargin=*]
\item Gauss-point evaluation through \texttt{SourceDriver(...)}
\item element integration through \texttt{ApplyXx4(...)} and
\texttt{Gauss2Node(...)}
\end{itemize}

\paragraph{Term 3: flux volume term $(F, \nabla \phi)_K$}

Still in \texttt{RuElemBlock(...)}:

\begin{itemize}[leftmargin=*]
\item \texttt{FluxDriver(...)} evaluates the physical flux
$F(u,q,o,w)$ at Gauss points
\item \texttt{ApplyXx4(...)} applies the metric/Jacobian transformation and
contracts with basis gradients
\item \texttt{Gauss2Node(...)} accumulates the result into the nodal residual
\end{itemize}

This is the discrete implementation of the LDG volume flux term.

So in Exasim:

\begin{itemize}[leftmargin=*]
\item \texttt{FluxDriver(...)} gives the PDE-specific flux
\item \texttt{ApplyXx4(...)} gives the mapped weak derivative action
\item \texttt{Gauss2Node(...)} completes the quadrature-based residual assembly
\end{itemize}

\subsection*{E. Implementation of the face term $\langle \hat{f}, \phi \rangle_{\partial K}$}

This is handled in \texttt{RuFaceBlock(...)} in
\texttt{backend/Discretization/uresidual.cpp}.

What the code does:

\begin{itemize}[leftmargin=*]
\item gathers \texttt{uh}, interior/exterior \texttt{udg}, and optional
\texttt{wdg} on faces
\item interpolates them to face Gauss points with \texttt{Node2Gauss(...)}
\item for interior faces:
\begin{itemize}[leftmargin=*]
\item calls \texttt{FhatDriver(...)}
\end{itemize}
\item for boundary faces:
\begin{itemize}[leftmargin=*]
\item calls \texttt{FbouDriver(...)}
\end{itemize}
\item multiplies by face Jacobian with \texttt{ApplyJacFhat(...)}
\item projects back with \texttt{Gauss2Node(...)}
\end{itemize}

Interpretation:

\begin{itemize}[leftmargin=*]
\item \texttt{FhatDriver(...)} is the interior numerical flux term of LDG
\item \texttt{FbouDriver(...)} is the boundary numerical flux term of LDG
\end{itemize}

This is exactly the implementation of $\langle \hat{f}, \phi \rangle_{\partial K}$.

\subsection*{F. How $\hat{f}$ is built}

The most important PDE-dependent LDG face term is in
\texttt{backend/Model/ModelDispatch/model\_drivers\_abi.cpp}.

\texttt{FhatDriver(...)} has two modes.

\paragraph{Mode 1: external/provider numerical flux}

If \texttt{common.extFhat == 1}, Exasim directly calls:

\begin{itemize}[leftmargin=*]
\item \texttt{abi.KokkosFhat(...)}
\end{itemize}

That means the model/provider supplies the full numerical trace flux.

\paragraph{Mode 2: Exasim-built LDG flux}

Otherwise, Exasim builds the numerical flux itself:

\begin{enumerate}[leftmargin=*]
\item compute left and right physical fluxes via \texttt{FluxDriver(...)}
\item average them with \texttt{AverageFlux(...)}
\item form the normal component with \texttt{AverageFluxDotNormal(...)}
\item add stabilization via:
\begin{itemize}[leftmargin=*]
\item \texttt{abi.KokkosStab(...)} if \texttt{common.extStab >= 1}
\item or built-in \texttt{AddStabilization1/2/3(...)}
\end{itemize}
\end{enumerate}

So the LDG numerical flux term is implemented as:

\begin{itemize}[leftmargin=*]
\item average normal physical flux
\item plus stabilization/jump penalty
\end{itemize}

That is the concrete code realization of the numerical flux in the weak face
integral.

\subsection*{G. Geometry and quadrature operators}

All volume and face weak terms rely on geometry computed in:

\begin{itemize}[leftmargin=*]
\item \texttt{backend/Discretization/geometry.cpp}
\end{itemize}

\texttt{ElemGeom(...)} computes:

\begin{itemize}[leftmargin=*]
\item physical Gauss-point coordinates \texttt{xg}
\item metric terms \texttt{Xx}
\item element Jacobian \texttt{jac}
\end{itemize}

\texttt{FaceGeom(...)} computes:

\begin{itemize}[leftmargin=*]
\item face coordinates
\item normals \texttt{nlg}
\item face Jacobians
\end{itemize}

These are consumed by:

\begin{itemize}[leftmargin=*]
\item \texttt{ApplyXx3(...)}
\item \texttt{ApplyXx4(...)}
\item \texttt{ApplyJacNormal(...)}
\item \texttt{ApplyJacFhat(...)}
\end{itemize}

So every weak-form term is implemented in the standard finite-element
sequence:

\begin{itemize}[leftmargin=*]
\item reference-to-physical geometry map
\item Gauss-point evaluation
\item metric/Jacobian scaling
\item quadrature projection back to nodal residuals
\end{itemize}

\subsection*{H. Optional auxiliary field $w$}

Exasim can also carry an auxiliary variable $w$ through \texttt{GetW(...)} in
\texttt{backend/Discretization/residual.cpp}.

This is not part of the core LDG $u$-$q$ structure, but it modifies source and
flux evaluations when \texttt{common.ncw > 0}.

Then:

\begin{itemize}[leftmargin=*]
\item \texttt{SourceDriver(...)} and \texttt{FluxDriver(...)} can depend on $w$
\item the Jacobian assembly also includes $w$-coupling corrections
\end{itemize}

For a basic LDG description, this is an extension layer, not the essential
structure.

\section*{4. Proposed implementation}

The clean way to understand Exasim's LDG implementation is this mapping from
weak form to code:

\begin{itemize}[leftmargin=*]
\item $q$ equation
\begin{itemize}[leftmargin=*]
\item $(q, v)_K$
\begin{itemize}[leftmargin=*]
\item local mass inversion in \texttt{GetQ(...)} via \texttt{ApplyMinv(...)}
\end{itemize}
\item $(u, \nabla \cdot v)_K$
\begin{itemize}[leftmargin=*]
\item \texttt{RqElemBlock(...)} in \texttt{backend/Discretization/qresidual.cpp}
\end{itemize}
\item $\langle \hat{u}, v \cdot n \rangle_{\partial K}$
\begin{itemize}[leftmargin=*]
\item \texttt{RqFaceBlock(...)} in \texttt{backend/Discretization/qresidual.cpp}
\end{itemize}
\end{itemize}
\item face trace construction
\begin{itemize}[leftmargin=*]
\item interior $\hat{u}$
\begin{itemize}[leftmargin=*]
\item copied from face $u$ in \texttt{GetUhat(...)}
\end{itemize}
\item boundary $\hat{u}$
\begin{itemize}[leftmargin=*]
\item \texttt{UbouDriver(...)} in
\texttt{backend/Model/ModelDispatch/model\_drivers\_abi.cpp}
\end{itemize}
\end{itemize}
\item $u$ equation
\begin{itemize}[leftmargin=*]
\item $(u_t, \phi)_K$
\begin{itemize}[leftmargin=*]
\item time-source assembly in \texttt{RuElemBlock(...)}
\end{itemize}
\item $(S, \phi)_K$
\begin{itemize}[leftmargin=*]
\item \texttt{SourceDriver(...)} in \texttt{RuElemBlock(...)}
\end{itemize}
\item $(F, \nabla \phi)_K$
\begin{itemize}[leftmargin=*]
\item \texttt{FluxDriver(...)} plus \texttt{ApplyXx4(...)} in
\texttt{RuElemBlock(...)}
\end{itemize}
\item $\langle \hat{f}, \phi \rangle_{\partial K}$
\begin{itemize}[leftmargin=*]
\item \texttt{FhatDriver(...)} or \texttt{FbouDriver(...)} in
\texttt{RuFaceBlock(...)}
\end{itemize}
\end{itemize}
\item geometry support
\begin{itemize}[leftmargin=*]
\item element metrics and Jacobians
\begin{itemize}[leftmargin=*]
\item \texttt{ElemGeom(...)}
\end{itemize}
\item face normals and Jacobians
\begin{itemize}[leftmargin=*]
\item \texttt{FaceGeom(...)}
\end{itemize}
\end{itemize}
\end{itemize}

\section*{5. Risks or numerical concerns}

\begin{itemize}[leftmargin=*]
\item The biggest source of confusion is sign convention. The textbook weak
form and the algebraic residual in Exasim may differ by an overall sign, but
the code comments in \texttt{qresidual.cpp} and \texttt{uresidual.cpp} make the
implemented structure clear.
\item \texttt{uh} in LDG is not an HDG global trace unknown in Exasim. It is a
face state used to construct numerical fluxes and boundary data.
\item The model-driver layer is critical: the mathematical LDG terms are
generic, but the actual PDE content is injected only through:
\begin{itemize}[leftmargin=*]
\item \texttt{FluxDriver(...)}
\item \texttt{SourceDriver(...)}
\item \texttt{FhatDriver(...)}
\item \texttt{FbouDriver(...)}
\item \texttt{UbouDriver(...)}
\item \texttt{TdfuncDriver(...)}
\end{itemize}
\item If \texttt{ncw > 0}, the implementation includes extra algebraic couplings
beyond the basic LDG $u$-$q$ structure.
\item For Jacobian-level analysis, the corresponding matrix assembly files such
as \texttt{backend/Discretization/qequation.cpp} and
\texttt{backend/Discretization/uequation.cpp} matter, but they are not the
primary residual-level LDG path.
\end{itemize}

\section*{6. Verification plan}

To verify this interpretation in the codebase:

\begin{enumerate}[leftmargin=*]
\item Start from \texttt{backend/Discretization/residual.cpp} and trace:
\begin{itemize}[leftmargin=*]
\item \texttt{GetUhat(...)}
\item \texttt{GetQ(...)}
\item \texttt{RuElem(...)}
\item \texttt{RuFace(...)}
\end{itemize}
\item For the weak $q$ equation, inspect:
\begin{itemize}[leftmargin=*]
\item \texttt{backend/Discretization/qresidual.cpp}
\end{itemize}
\item For the weak $u$ equation, inspect:
\begin{itemize}[leftmargin=*]
\item \texttt{backend/Discretization/uresidual.cpp}
\end{itemize}
\item For the numerical flux and boundary terms, inspect:
\begin{itemize}[leftmargin=*]
\item \texttt{backend/Model/ModelDispatch/model\_drivers\_abi.cpp}
\end{itemize}
\item For a concrete PDE, test the interpretation on a simple example such as
Poisson in first-order form and confirm:
\begin{itemize}[leftmargin=*]
\item $q$ is recovered locally
\item $u$ residual receives volume flux, source, and face flux terms
\item boundary data enters through \texttt{UbouDriver(...)} and
\texttt{FbouDriver(...)}
\end{itemize}
\end{enumerate}

\end{document}
